% !TeX root = ../se200_referential_regimes.tex

% =============================================================================
\section{Determinacy}
\label{sec:determinacy}
% =============================================================================

The lower-bound result requires a fixed assignment
from core carrier-basis pairs to identity regimes.
This section records that assignment and
only addresses the carrier-basis pairs used in the lower-bound construction.

Once a core carrier-basis pair is fixed,
the regime assignment is determined.
Which identity basis a real-world referent ought to use requires modeling judgment.

% -----------------------------------------------------------------------------
\subsection{Core Carrier-Basis Pairs}
\label{subsec:core_carrier_basis_pairs}
% -----------------------------------------------------------------------------

\begin{definition}[Core carrier-basis pair]
\label{def:core_carrier_basis_pair}
A \emph{core carrier-basis pair} is one of the carrier kind and identity basis
combinations used in the lower-bound construction.
The core pairs are:
\[
\begin{array}{rcl}
\text{obligation-bearer with responsibility-bearing identity} &\mapsto& \OBL,\\
\text{occurrence with realization-and-provenance identity} &\mapsto& \OCC,\\
\text{record with descriptive-record identity} &\mapsto& \REC,\\
\text{plain referent with locus-fixed identity} &\mapsto& \LOC,\\
\text{plain referent with object-fixed identity} &\mapsto& \OBJ,\\
\text{scope with extension-fixed identity} &\mapsto& \SCOPEE,\\
\text{scope with structure-fixed identity} &\mapsto& \SCOPES,\\
\text{rule with content-fixed identity} &\mapsto& \RULEC,\\
\text{rule with structure-fixed identity} &\mapsto& \RULES.
\end{array}
\]
\end{definition}

The first three pairs contribute one regime each.
The remaining six pairs arise from the witnessed split cases proved in Section~\ref{sec:algebra}.

\begin{definition}[Core regime assignment]
\label{def:core_regime_assignment}
The \emph{core regime assignment} maps each core carrier-basis pair
to the identity regime displayed in Definition~\ref{def:core_carrier_basis_pair}.
\end{definition}

% -----------------------------------------------------------------------------
\subsection{Existence}
\label{subsec:determinacy_existence}
% -----------------------------------------------------------------------------

\begin{lemma}[Existence for the core]
\label{se200.lem.CoreExistence}
Every core carrier-basis pair is assigned to at least one of the nine core regimes.
\end{lemma}

\begin{proof}
\sloppy
The three one-basis core pairs are assigned by construction:
obligation-bearer with responsibility-bearing identity maps to $\OBL$,
occurrence with realization-and-provenance identity maps to $\OCC$,
and record with descriptive-record identity maps to $\REC$.

The remaining core pairs are assigned by the split results of Section~\ref{sec:algebra}.
A plain referent with locus-fixed identity maps to $\LOC$,
and a plain referent with object-fixed identity maps to $\OBJ$,
by Proposition~\ref{se200.prop.PlainReferentUnderdetermination}.
A scope with extension-fixed identity maps to $\SCOPEE$,
and a scope with structure-fixed identity maps to $\SCOPES$,
by Proposition~\ref{se200.prop.ScopeUnderdetermination}.
A rule with content-fixed identity maps to $\RULEC$,
and a rule with structure-fixed identity maps to $\RULES$,
by Proposition~\ref{se200.prop.RuleUnderdetermination}.
Therefore every core carrier-basis pair is assigned to at least one of the nine core regimes.
\end{proof}

% -----------------------------------------------------------------------------
\subsection{Uniqueness}
\label{subsec:determinacy_uniqueness}
% -----------------------------------------------------------------------------

\begin{lemma}[Uniqueness for the core]
\label{se200.lem.CoreUniqueness}
No core carrier-basis pair is assigned to more than one of the nine core regimes.
\end{lemma}

\begin{proof}
Each core carrier-basis pair contains a fixed carrier kind and a fixed identity basis.
The core regime assignment maps that pair
to the regime that preserves the declared basis
under the transformation behavior used in the lower-bound construction.

The three one-basis core pairs map to $\OBL$, $\OCC$, and $\REC$ respectively.
For the split cases, the declared basis selects the side of the split:
locus-fixed plain referent maps to $\LOC$ and not $\OBJ$;
object-fixed plain referent maps to $\OBJ$ and not $\LOC$;
extension-fixed scope maps to $\SCOPEE$ and not $\SCOPES$;
structure-fixed scope maps to $\SCOPES$ and not $\SCOPEE$;
content-fixed rule maps to $\RULEC$ and not $\RULES$;
and structure-fixed rule maps to $\RULES$ and not $\RULEC$.
Thus the assignment is functional on the core pairs.
Therefore no core carrier-basis pair is assigned to more than one core regime.
\end{proof}

% -----------------------------------------------------------------------------
\subsection{Determinacy Theorem}
\label{subsec:determinacy_theorem}
% -----------------------------------------------------------------------------

\begin{theorem}[Core regime assignment]
\label{se200.thm.CoreRegimeAssignment}
Every core carrier-basis pair is assigned to exactly one of the nine core identity regimes.
\end{theorem}

\begin{proof}
\sloppy
Existence follows from Lemma~\ref{se200.lem.CoreExistence}.
Uniqueness follows from Lemma~\ref{se200.lem.CoreUniqueness}.
Therefore every core carrier-basis pair is assigned to exactly one of the nine core regimes.
\end{proof}

This determines the core assignment.
It does not decide whether a modeler has chosen the correct identity basis
for a particular real-world referent.
It says only that the carrier-basis pairs used in the lower-bound construction
have fixed regime assignments.

% -----------------------------------------------------------------------------
\subsection{Failure of Basis Declaration}
\label{subsec:basis_declaration_failure}
% -----------------------------------------------------------------------------

For the witnessed split cases, omission of the basis is not harmless.
A scope without an extension-fixed or structure-fixed basis does not specify
whether decomposition preserves covered cases or internal structure.
A rule without a content-fixed or structure-fixed basis does not specify
whether refinement preserves rule content or textual structure.
A plain referent without a locus-fixed or object-fixed basis does not specify
whether forking preserves location or object identity.
A declared identity basis is required for deterministic assignment.

This observation anticipates the conformance analysis
of the companion \emph{Accountable Records} paper.
In a deployed record system, basis declaration becomes a schema obligation.
In this paper, it marks the boundary of the determinacy result:
regime assignment is fixed for the carrier-basis pairs included
in the core.
